\documentclass[journal,onecolumn,11pt]{IEEEtran}

\usepackage{amsmath,amssymb,amsfonts,amsthm}
\usepackage{algorithmic}
\usepackage{algorithm}
\usepackage{array}
\usepackage{booktabs}
\usepackage{cite}
\usepackage{graphicx}
\usepackage{textcomp}
\usepackage{url}
\usepackage{hyperref}
\usepackage{multirow}

\newtheorem{theorem}{Theorem}
\newtheorem{lemma}{Lemma}
\newtheorem{definition}{Definition}
\newtheorem{corollary}{Corollary}
\newtheorem{remark}{Remark}

\begin{document}

\title{Dumb Ways to Solve Subset Sum: Exact Single-Threaded Solving at $N = 100$, $T \approx 1\text{T}$, Density $\approx 2.51$, and Under $50$\,MB RAM on a Low-Cost Laptop}

\author{M. Diki Fahriza\\
Independent Researcher\\
\texttt{mdikifahriza3@gmail.com}}

\maketitle

\begin{abstract}
The Subset Sum Problem (SSP) is a foundational NP-complete problem with applications in cryptanalysis, integer programming, and combinatorial optimization. Practical exact solvers face three recurring bottlenecks: symmetric Meet-in-the-Middle (MITM) requires $\mathcal{O}(2^{N/2})$ memory and causes severe CPU cache thrashing once $N \gtrsim 50$; classical dynamic programming (DP) requires $\mathcal{O}(T)$ memory and collapses on trillion-scale targets ($T \ge 10^{12}$); and extracting many distinct witnesses from a branch-and-bound tree forces expensive re-traversal. This paper presents \textsc{Dumb SSP Solver}, a native C++20 exact solver architecture that integrates four concrete engineering mechanisms: (1) an \emph{asymmetric tail-table memoization engine} over the $m$ numerically smallest active elements ($m = N-1$ for $N \le 20$, and $m = 20$ bounded by memory for $N > 20$), restricting table footprint to $16.00$\,MB to reside entirely within CPU Level-3 (L3) cache while pruning the last $m$ recursion levels via a single binary search per leaf; (2) a $64$-bit word-parallel bitset DP engine restricted to small effective targets ($T \le 1.5\times10^{7}$) and single-witness queries (\texttt{FindOne}/\texttt{DecisionOnly}); (3) a bounded-degree \emph{Zero-Sum Swap Extractor} that harvests additional witnesses by exploring symmetric and asymmetric element substitutions up to degree $p \le 4$ in fixed-degree polynomial time $\mathcal{O}(k^4\log k + (N-k)^4\log k)$; and (4) an independent, isolated $\mathcal{O}(k)$ black-box verifier (L7) that validates index bounds, element uniqueness, and exact sum equality outside the search engine. All reported timings are derived from a fully reproducible single-threaded self-test benchmark spanning $N=5$ to $N=100$ with trillion-scale targets, executed on commodity hardware and independently verifiable via \texttt{dumbsspCli.exe}.
\end{abstract}

\begin{IEEEkeywords}
Subset Sum Problem, Meet-in-the-Middle, Algorithm Engineering, Cache Locality, Bitset Dynamic Programming, Local Search Neighborhood, Zero-Sum Perturbation, Parameterized Complexity, Reproducibility.
\end{IEEEkeywords}

\section{Introduction}
\label{sec:introduction}

\IEEEPARstart{T}{he} \textsc{Subset Sum Problem} (SSP) is one of the most celebrated and fundamental problems in theoretical computer science, operations research, and combinatorial optimization. Formally, given an ordered multiset of $N$ positive integers $A = \{a_1, a_2, \dots, a_N\} \subset \mathbb{Z}^+$ and a prescribed target integer $T \in \mathbb{Z}^+$, the decision problem asks whether there exists a subset of indices $S \subseteq \{1, 2, \dots, N\}$ such that the summation of the selected elements equals exactly $T$:
\begin{equation}
    \sum_{i \in S} a_i = T.
    \label{eq:ssp_core}
\end{equation}

Three formulations are considered throughout this paper and are all directly supported by the implementation as distinct \texttt{SolveMode} values:
\begin{enumerate}
    \item \textbf{Decision Variant (\texttt{DecisionOnly}):} determine the existence of at least one satisfying subset, $\Phi(A, T) \in \{\text{TRUE}, \text{FALSE}\}$.
    \item \textbf{Search Variant (\texttt{FindOne}/\texttt{FindAll}):} reconstruct an explicit binary witness $x = (x_1, \dots, x_N) \in \{0, 1\}^N$ with $\sum_{i=1}^N a_i x_i = T$, either a single witness or as many as the configured cap allows.
    \item \textbf{Counting Variant (\texttt{CountAll}):} compute the cardinality
    \begin{equation}
        \mathcal{N}(A, T) = \left| \left\{ S \subseteq \{1, 2, \dots, N\} : \sum_{i \in S} a_i = T \right\} \right|
    \end{equation}
    over the \emph{active} elements of $A$ (Section~\ref{sec:preliminaries}); see Remark~\ref{rem:zero-count} for a caveat regarding zero-valued input elements.
\end{enumerate}

\subsection{Complexity-Theoretic Landscape}
Since Karp's seminal 1972 treatise on combinatorial reducibility \cite{karp1972reducibility}, SSP has maintained a central role in computational complexity theory. As one of Karp's 21 NP-complete problems, SSP is intractable in the worst case under $\text{P} \neq \text{NP}$, and under the Strong Exponential Time Hypothesis exact algorithms for arbitrary instances are conjectured to require $\mathcal{O}^*(2^{(1-\varepsilon)N})$ time.

SSP is nevertheless only \textit{weakly} NP-complete: when $T$ and the $a_i$ are bounded polynomially in $N$, Bellman's dynamic programming solves SSP in pseudo-polynomial time $\mathcal{O}(N \cdot T)$ \cite{pisinger1997minimal}. Near-linear pseudo-polynomial algorithms running in $\widetilde{\mathcal{O}}(N + T)$ time have since been developed via color coding, generating functions, and FFT convolution (Bringmann \cite{bringmann2017near}, Jin and Wu \cite{jin2019simple}, Chen et al. \cite{chen2024improved}).

\subsection{Practical Bottlenecks in Existing Solvers}
\begin{enumerate}
    \item \textbf{The MITM Memory Wall:} Horowitz--Sahni \cite{horowitz1974computing} partitions $A$ into two halves of size $\lfloor N/2 \rfloor$ and $\lceil N/2 \rceil$, generating, sorting, and matching $2^{N/2}$ partial sums in $\mathcal{O}(2^{N/2})$ time and space. For $N \ge 50$, $2^{25} \approx 3.36 \times 10^7$ stored entries already exceed typical CPU L3 cache sizes ($16$--$64$\,MB on commodity desktop/server parts), causing cache thrashing and TLB pressure. Schroeppel and Shamir \cite{schroeppel1981t} reduce space to $\mathcal{O}(2^{N/4})$ at the cost of priority-queue overhead.
    \item \textbf{The Target-Magnitude Wall:} Bellman DP and near-linear algorithms are excellent while $T$ stays within tens of millions, but a plain DP table of size $\mathcal{O}(T)$ is simply infeasible once $T$ reaches $10^{12}$.
    \item \textbf{The Multi-Solution Extraction Penalty:} obtaining $K \gg 1$ distinct witnesses by naive re-traversal of a branch-and-bound tree costs $\mathcal{O}(K \cdot 2^N)$ in the worst case.
\end{enumerate}

\subsection{Contributions and Scope}
This paper documents the architecture, implementation, and empirical behavior of \textsc{Dumb SSP Solver}, implemented in native C++20 (\texttt{dumbsspCore.hpp}, with a synchronous console front-end \texttt{dumbsspCli.cpp} and an interactive Win32 graphical interface \texttt{dumbsspGui.cpp}):
\begin{itemize}
    \item \textbf{A deterministic multi-stage pre-check pipeline (L0--L2)} that performs domain filtering, a GCD obstruction sieve, and a cardinality-window feasibility test in polynomial time prior to exponential search.
    \item \textbf{A memory-bounded asymmetric tail-table (L4)} over the $m$ numerically smallest active elements, maintaining a compact $16.00$\,MB footprint that resides in L3 CPU cache and accelerates leaf evaluation.
    \item \textbf{A word-parallel bitset DP engine (L3)}, threshold-gated to small target regimes ($T \le 1.5\times10^7$) for high-throughput single-witness resolution.
    \item \textbf{A bounded-degree Zero-Sum Swap Extractor (L8)}, establishing a fixed-degree polynomial framework for multi-solution harvesting up to degree $p \le 4$.
    \item \textbf{A fully reproducible benchmark suite}: providing standalone per-instance datasets and exact CLI invocations from $N=5$ to $N=100$ with trillion-scale targets, archived in \texttt{benchmark.txt}.
\end{itemize}

\subsection{Paper Organization}
Section~\ref{sec:preliminaries} gives the formal model, the Complement Invariance Reduction theorem, and the (diagnostic-only) density metric. Section~\ref{sec:architecture} details the L0--L2 pre-check pipeline and the exact router logic (Algorithm~\ref{alg:router}), reproduced from \texttt{AdaptiveStrategySelector::select}. Section~\ref{sec:hardware_engineering} covers the bitset DP, the tail-table engine, and the pruning oracle. Section~\ref{sec:swap_manifold} gives the Zero-Sum Swap theory and its corrected complexity. Section~\ref{sec:experiments} presents the reproducible self-test benchmark and reproduction instructions. Section~\ref{sec:ui_verification} discusses the native GUI, a known reliability caveat, and the L7 verifier. Section~\ref{sec:conclusion} concludes.

\section{Preliminaries \& Complexity-Theoretic Taxonomy}
\label{sec:preliminaries}

\subsection{Formal Mathematical Formulation}
Let $A_{\text{raw}}$ be the $N$-element input multiset as supplied by the user. The solver first constructs the \emph{active} multiset
\begin{equation}
    A = \{\, a \in A_{\text{raw}} : a \neq 0 \ \wedge\ a \le T \,\},
    \label{eq:active_set}
\end{equation}
discarding zero-valued elements (they never change a subset sum) and elements strictly larger than $T$ (a positive element larger than $T$ can never belong to any exact solution, since all elements are positive). Elements of $A$ are then sorted into non-increasing order, $a_1 \ge a_2 \ge \dots \ge a_n > 0$ with $n = |A| \le N$. We write $\sigma(A) = \sum_{i=1}^{n} a_i$ for the aggregate sum of the \emph{active} multiset only (not of $A_{\text{raw}}$).

\begin{lemma}[Soundness of Domain Filtering]
\label{lem:filter}
For any $S_{\text{raw}} \subseteq A_{\text{raw}}$ with $\sum S_{\text{raw}} = T$: (i) $S_{\text{raw}}$ contains no zero element by assumption of positivity of a witness of interest, and any zero elements may be freely added to or removed from $S_{\text{raw}}$ without changing the sum; (ii) $S_{\text{raw}}$ contains no element $a > T$, because $a \le \sum S_{\text{raw}} = T$ is required for every $a \in S_{\text{raw}}$ when all elements are positive.
\end{lemma}
\begin{proof}
(ii) follows because all elements of $A_{\text{raw}}$ are non-negative integers; if $a \in S_{\text{raw}}$ and $a > T$, then $\sum S_{\text{raw}} \ge a > T$, contradicting $\sum S_{\text{raw}} = T$. (i) is immediate since adding $0$ to a sum leaves it unchanged.
\end{proof}

A candidate solution is represented as $x = (x_1, \dots, x_n) \in \{0,1\}^n$ over the active set, with $f(x) = \sum_{i=1}^n a_i x_i$.

\subsection{The Complement Invariance Reduction Theorem}
\begin{theorem}[Complement Invariance Reduction]
\label{thm:complement}
For an active instance $(A,T)$ with $T > \tfrac{1}{2}\sigma(A)$ and $T \le \sigma(A)$, finding $x \in \{0,1\}^n$ with $\sum_{i=1}^n a_i x_i = T$ is equivalent to finding $x' = \mathbf{1}-x$ with
\begin{equation}
    \sum_{i=1}^n a_i x'_i = T' = \sigma(A) - T, \qquad T' < \tfrac{1}{2}\sigma(A).
\end{equation}
Given a certified witness for $T'$, $x$ is recovered in $\mathcal{O}(n)$ time via $x_i = 1 - x'_i$.
\end{theorem}
\begin{proof}
Let $x'_i = 1-x_i$. Then $\sum_i a_i x'_i = \sum_i a_i - \sum_i a_i x_i = \sigma(A) - T = T'$. Since $T > \tfrac12\sigma(A)$, we get $T' = \sigma(A)-T < \tfrac12\sigma(A)$. The converse direction is symmetric, giving a bijection between the two solution sets.
\end{proof}
This is exactly the transformation performed by \texttt{Instance::normalize()}: \texttt{effective\_target} is set to $\sigma(A)-T$ and \texttt{complement\_applied} is raised whenever $T > \lfloor \sigma(A)/2 \rfloor$ and $\sigma(A) \ge T$; every returned witness is inverted back against $A$ (not $A_{\text{raw}}$) before the final L7 verification pass.

\subsection{Diagnostic-Only Information Density}
The implementation also computes an information-density score,
\begin{equation}
    d = \frac{n}{\log_2(a_1 + 1)}, \qquad n = |A|,\ a_1 = \max A,
    \label{eq:density_formula}
\end{equation}
and a boolean \texttt{strong\_structure} flag from the cardinality window and GCD (defined in Section~\ref{sec:architecture}). Both quantities are reported to the user in the metadata panel for interpretability.

\begin{remark}[Diagnostic Role of Information Density]
\label{rem:density-not-routing}
Information density $d$ and the structural flag are evaluated as informative metrics in the telemetry panel to characterize theoretical instance difficulty according to classical phase-transition literature \cite{lagarias1985solving, coster1992improved}. Strategy dispatch is governed deterministically by memory bounds and target thresholds as formulated in Algorithm~\ref{alg:router}.
\end{remark}

For context, the classical regimes are: \textbf{low density} ($d < 0.9408$), where Lagarias and Odlyzko \cite{lagarias1985solving} and Coster et al. \cite{coster1992improved} demonstrated that random instances almost always reduce in polynomial time to finding a short vector in an $(N{+}1)$-dimensional lattice via LLL; \textbf{critical density} ($d \approx 1$), where exact search is asymptotically hardest; and \textbf{high density} ($d>1.2$), where pigeonhole effects create many redundant solutions and dynamic programming is effective.

\subsection{Complexity-Theoretic Taxonomy of Solver Layers}
Table~\ref{tab:taxonomy} classifies every module exactly as implemented; the space/time entries are derived directly from the code paths cited in the ``Mechanism'' column, not from asymptotically idealized versions of the same idea.

\begin{table*}[t]
\centering
\caption{Complexity-Theoretic Classification of the Solver Architecture (as implemented)}
\label{tab:taxonomy}
\renewcommand{\arraystretch}{1.3}
\footnotesize
\resizebox{\textwidth}{!}{%
\begin{tabular}{lllll}
\hline
\textbf{Layer / Module} & \textbf{Class} & \textbf{Time} & \textbf{Space} & \textbf{Mechanism} \\
\hline
L0: Filter \& Normalizer & P & $\mathcal{O}(N)$ & $\mathcal{O}(n)$ & Drop zeros/oversized elements, apply Thm.~\ref{thm:complement} \\
L1: GCD Sieve & P & $\mathcal{O}(n \log(\min A_i))$ & $\mathcal{O}(1)$ & $T \not\equiv 0 \!\!\pmod{\gcd(A)} \Rightarrow$ UNSAT \\
L2: Cardinality Profiler & P & $\mathcal{O}(n \log n)$ & $\mathcal{O}(n)$ & Sort, suffix sums, window $[k_{\min},k_{\max}]$ \\
L3: Word-Parallel Bitset DP & Pseudo-P. & $\mathcal{O}(n\, T/64 + T)$ & $\mathcal{O}(T)$ & 64-bit shift DP; only if $T \le 1.5\times10^7$, mode $\in\{$FindOne, DecisionOnly$\}$ \\
L4: Asymmetric Tail-Table & Param.\ Exp. & $\mathcal{O}(m2^{m} + 2^{\,n-m})$ & $\mathcal{O}(2^{m})$ & Precomputed table over the $m$ smallest active elements (Sec.~\ref{sec:hardware_engineering}) \\
L5a: Suffix Bound Prune & P & $\mathcal{O}(1)$/node & $\mathcal{O}(n)$ & $\text{rem} > \text{suffix}[i] \Rightarrow$ prune \\
L5b: Local Cardinality Oracle & P & $\mathcal{O}(n{-}i)$/node & $\mathcal{O}(n)$ & Binary search + local $[k_{\min},k_{\max}]$ re-check (Sec.~\ref{sec:hardware_engineering}) \\
L6: Cooperative Cancellation & P & $\mathcal{O}(1)$ amort.\ & $\mathcal{O}(1)$ & One atomic flag, polled every $4096$ nodes; \emph{not} multi-threaded search \\
L7: Independent Verifier & P & $\mathcal{O}(k)$ & $\mathcal{O}(k)$ & Isolated re-check: bounds, uniqueness, $\sum_{j\in S} a_j \stackrel{?}{=} T$ \\
L8: Bounded-Degree Swap Extr. & P (fixed deg.) & $\mathcal{O}(k^4\log k + (n{-}k)^4\log k)$ & $\mathcal{O}(k^4 + (n{-}k)^4)$ interm. & Exchange $p\le 4$ elements per side (Sec.~\ref{sec:swap_manifold}) \\
\hline
\end{tabular}%
}
\end{table*}

\section{Pre-Check Pipeline and the Adaptive Router}
\label{sec:architecture}

The solver is not a single monolithic algorithm; three deterministic pre-check layers run first, and only then does a router choose between two exponential/pseudo-polynomial search engines.

\subsection{Pre-Processing, Sieve, and Profiling Layers (L0--L2)}
\begin{enumerate}
    \item \textbf{L0 (Filter \& Normalizer):} builds $A$ from $A_{\text{raw}}$ per Eq.~\eqref{eq:active_set} and Lemma~\ref{lem:filter}, computes $\sigma(A)$, and applies Theorem~\ref{thm:complement} when $T > \tfrac12\sigma(A)$.
    \item \textbf{L1 (GCD Sieve):} computes $g=\gcd(a_1,\dots,a_n)$ in $\mathcal{O}(n\log(\min A_i))$ time.
    \begin{lemma}[Modular Obstruction]
    If $T \not\equiv 0 \pmod g$ then no $S \subseteq A$ sums to $T$.
    \end{lemma}
    \begin{proof}
    $\sum_{i\in S} a_i = g\sum_{i\in S}(a_i/g)$ is a multiple of $g$ for any $S$.
    \end{proof}
    The implementation also performs a redundant $\mathcal{O}(1)$ fast-path check (all active elements even, effective target odd $\Rightarrow$ UNSAT), which is a special case already implied by the GCD test whenever $g>1$; it exists purely to short-circuit that common case without waiting for the general GCD computation to be interpreted.
    \item \textbf{L2 (Cardinality Profiler):} sorts $A$ descending, computes the suffix sums $\text{suffix}[i]=\sum_{j=i}^{n} a_j$, and the cardinality window
    \begin{align}
        k_{\min} &= \min\Big\{k : \textstyle\sum_{i=1}^k a_i \ge T_{\text{eff}}\Big\}, \\
        k_{\max} &= \max\Big\{k : \textstyle\sum_{i=n-k+1}^n a_i \le T_{\text{eff}}\Big\},
    \end{align}
    i.e.\ $k_{\min}$ is the fewest \emph{largest} elements that can reach $T_{\text{eff}}$, and $k_{\max}$ the most \emph{smallest} elements that do not overshoot it. If $T_{\text{eff}}=0$, $k_{\min}=k_{\max}=0$ by convention (the empty set). If $k_{\min}>k_{\max}$, the instance is certified UNSAT.
\end{enumerate}

\subsection{Adaptive Router Decision Logic}
Algorithm~\ref{alg:router} is reproduced from \texttt{AdaptiveStrategySelector::select}, including its exact numeric thresholds.

\begin{algorithm}[h]
\caption{Adaptive Structural Router Dispatch (verbatim thresholds)}
\label{alg:router}
\begin{algorithmic}[1]
\REQUIRE Active multiset $A$, effective target $T_{\text{eff}}$, memory budget $M_{\text{limit}}$ (MB), mode $\mathcal{M}$
\ENSURE Strategy $\mathcal{S} \in \{\text{TrivialPreCheck}, \text{BitsetDP}, \text{HybridTailTable}\}$
\IF{$T_{\text{eff}} = 0$}
    \RETURN $\mathcal{S}\leftarrow$ TrivialPreCheck (empty-set solution)
\ENDIF
\IF{$A=\emptyset$ \OR $T_{\text{eff}} > \sigma(A)$}
    \RETURN $\mathcal{S}\leftarrow$ TrivialPreCheck (UNSAT: exceeds sum)
\ENDIF
\IF{$g>1$ \AND $T_{\text{eff}} \bmod g \neq 0$}
    \RETURN $\mathcal{S}\leftarrow$ TrivialPreCheck (UNSAT: GCD sieve)
\ENDIF
\IF{all $a_i$ even \AND $T_{\text{eff}}$ odd}
    \RETURN $\mathcal{S}\leftarrow$ TrivialPreCheck (UNSAT: parity fast-path)
\ENDIF
\IF{$k_{\min} > k_{\max}$}
    \RETURN $\mathcal{S}\leftarrow$ TrivialPreCheck (UNSAT: empty cardinality window)
\ENDIF
\STATE $\text{Mem}_{\text{DP}} \leftarrow \big\lceil (T_{\text{eff}}+1)/64 \big\rceil \times 8 + (T_{\text{eff}}+1)\times 4$ \ \ (bytes; treated as $\infty$ if $T_{\text{eff}} > 5\times10^7$)
\IF{$T_{\text{eff}} \le 1.5\times 10^{7}$ \AND $\text{Mem}_{\text{DP}} < \tfrac12 M_{\text{limit}}\!\cdot\!2^{20}$ \AND $\mathcal{M} \in \{\text{FindOne}, \text{DecisionOnly}\}$}
    \RETURN $\mathcal{S}\leftarrow$ BitsetDP \ (Layer L3)
\ELSE
    \RETURN $\mathcal{S}\leftarrow$ HybridTailTable \ (Layers L4--L5)
\ENDIF
\end{algorithmic}
\end{algorithm}

Note in particular that \texttt{FindAll} and \texttt{CountAll} are \emph{never} routed to the bitset engine regardless of how small $T_{\text{eff}}$ is: only \texttt{FindOne} and \texttt{DecisionOnly} qualify. This differs from a superficially similar ``mode $\ne$ CountAll'' condition that would incorrectly also admit \texttt{FindAll}.

\subsection{Supported Exact Operational Modes}
\begin{itemize}
    \item \textbf{FindOne / DecisionOnly:} the search returns as soon as a leaf match is recorded, via the single \texttt{std::atomic<bool> solution\_found} flag checked at each recursive call.
    \item \textbf{FindAll:} after an initial witness is found, the (non-exhaustive, default) path hands off to the L8 Zero-Sum Swap Extractor (Section~\ref{sec:swap_manifold}) to harvest additional witnesses up to \texttt{max\_solutions} (default $5000$); an \texttt{exhaustive\_find\_all} flag instead lets the L4 depth-first search itself keep enumerating until the same cap or a wall-clock timeout is reached.
    \item \textbf{CountAll:} increments a $128$-bit counter (\texttt{u128 solution\_count}) for every match; after the first witness is captured (needed to populate \texttt{sample\_solution} for display), subsequent matches skip witness-vector allocation entirely.
\end{itemize}

\begin{remark}[Combinatorial Scaling for Zero-Valued Elements]
\label{rem:zero-count}
Zero-valued elements ($a_i = 0$) are filtered during domain reduction (L0) to preserve search tree efficiency. In counting and exhaustive enumeration modes, the total witness count over the original multiset algebraically corresponds to $\mathcal{N}(A_{\text{raw}}, T) = 2^Z \cdot \mathcal{N}(A, T)$, where $Z$ is the count of zero elements tracked in the diagnostic metadata.
\end{remark}

\section{Hardware-Conscious Algorithm Engineering}
\label{sec:hardware_engineering}

\subsection{Word-Parallel 64-Bit Bitset DP (L3)}
When routed to L3, reachability is stored as $\text{DP}[0..W{-}1]$, a word array of $W=\lceil (T_{\text{eff}}+1)/64\rceil$ 64-bit words, with $\text{DP}[0]$ initialized to bit $0$ set. For each active element $v$, inclusion corresponds to a multi-word left shift by $v$ bits:
\begin{equation}
    w_{\text{shift}} = \lfloor v/64\rfloor, \qquad b_{\text{rem}} = v \bmod 64,
\end{equation}
\begin{equation}
    \text{DP}'[w] = \text{DP}[w] \;\lor\; \big(\text{DP}[w-w_{\text{shift}}] \ll b_{\text{rem}}\big) \;\lor\; \big(\text{DP}[w-w_{\text{shift}}-1] \gg (64-b_{\text{rem}})\big).
    \label{eq:bitset_transition}
\end{equation}
This is exactly the update in \texttt{solve\_bitset}, processing $64$ reachability states per shifted word. A parallel \texttt{parent[t]} array of $T_{\text{eff}}{+}1$ signed 32-bit integers records, for every newly-reachable sum $t$, the index of the element that first reached it, enabling $\mathcal{O}(k)$ witness reconstruction by repeated subtraction. Because \texttt{parent} occupies $4(T_{\text{eff}}{+}1)$ bytes -- strictly larger than the $8\lceil (T_{\text{eff}}{+}1)/64\rceil$ bytes used by the bit-packed \texttt{DP} array -- the dominant space cost of L3 is $\mathcal{O}(T_{\text{eff}})$, not $\mathcal{O}(T_{\text{eff}}/64)$; we correct this in Table~\ref{tab:taxonomy}. Total time is $\mathcal{O}(n\,T_{\text{eff}}/64)$ for the word-level shifts plus $\mathcal{O}(T_{\text{eff}})$ amortized across the run for parent-array bookkeeping (each of the $T_{\text{eff}}{+}1$ states is claimed at most once).

\subsection{Asymmetric Tail-Table Memoization in CPU L3 Cache (L4)}
Symmetric MITM \cite{horowitz1974computing} splits $n$ elements $50{:}50$; for $n=50$ each half already needs $2^{25}\approx 3.36\times10^7$ entries. Layer L4 instead precomputes a table over only the $m$ \emph{numerically smallest} active elements, where $m$ is chosen by \texttt{HybridTailTableEngine::choose\_m} as:
\begin{equation}
    m(n, M_{\text{limit}}) =
    \begin{cases}
        0, & n \le 1,\\[2pt]
        n-1, & 2 \le n \le 20,\\[2pt]
        \displaystyle \max\Big(12,\ \max\{\, m' \le 20 : 16\cdot 2^{m'} \le \tfrac{M_{\text{limit}}\cdot 2^{20}}{4} \,\}\Big), & n > 20.
    \end{cases}
    \label{eq:choose_m}
\end{equation}
At the default budget $M_{\text{limit}}=4096$\,MB used by the CLI and GUI, the memory constraint in the third case never binds ($16\cdot 2^{20}=16$\,MB $\ll 1$\,GB), so $m=20$ for every $n>20$ tested in Section~\ref{sec:experiments}. For small instances ($n\le 20$), the tail table spans $m=n-1$ elements, reducing search to a single binary-search lookup over the entire active tail.

The table itself, $\text{Tail} = \{a_{n-m+1},\dots,a_n\}$, has exactly $2^m$ entries, each an $(\texttt{u64 sum}, \texttt{u32 mask})$ pair that occupies $16$ bytes under standard 8-byte alignment (8-byte \texttt{sum} forces the struct to a 16-byte stride). At $m=20$:
\begin{equation}
    \text{Memory}(\text{Table}_{m=20}) = 2^{20}\times 16\text{ bytes} = 16{,}777{,}216\text{ bytes} = 16.00\text{ MB},
\end{equation}
which fits entirely inside the L3 cache of contemporary desktop/server CPUs (typically $16$--$64$\,MB). The table is populated by evaluating partial sums across all $2^m$ bit combinations in $\mathcal{O}(m\,2^m)$ time, followed by an $\mathcal{O}(2^m\log 2^m)$ sort to facilitate binary search matching.

With \texttt{cutoff}~$=n-m$, the remaining $n-m$ (largest) elements are explored by a depth-first search over indices $0,\dots,\text{cutoff}-1$. Whenever the search reaches depth $i=\text{cutoff}$ with residual target $\text{rem}$, it queries the sorted table via \texttt{std::equal\_range} in $\mathcal{O}(\log 2^m)=\mathcal{O}(m)$ comparisons instead of continuing the recursion, which is what prunes (for the $n>20$, $m=20$ regime) the last $20$ recursion levels per branch.

\subsection{Suffix and Cardinality Pruning During DFS (L5)}
Two distinct pruning checks run at every DFS node $(i,\text{rem})$, and the implementation counts them separately (\texttt{states\_pruned} vs.\ \texttt{oracle\_calls}/\texttt{oracle\_pruned}):
\begin{itemize}
    \item \textbf{L5a (Suffix Bound, $\mathcal{O}(1)$):} if $\text{rem} > \text{suffix}[i] = \sum_{j=i}^{n} a_j$, the whole subtree is pruned immediately.
    \item \textbf{L5b (Local Cardinality Oracle, \texttt{is\_cardinality\_feasible}):} a strictly stronger, more expensive necessary condition. It (i) binary-searches for the largest remaining element $v\le \text{rem}$ in $\mathcal{O}(\log(n-i))$; (ii) rejects if fewer than $\lceil \text{rem}/v\rceil$ elements remain available; and (iii) rebuilds a \emph{local} cardinality window $[k_{\min}^{(i)}, k_{\max}^{(i)}]$ over indices $i,\dots,n-1$ by a linear scan, rejecting if it is empty. Step (iii) is $\mathcal{O}(n-i)$ in the worst case, so the oracle as a whole costs $\mathcal{O}(n-i) = \mathcal{O}(N)$ per call in the worst case, not $\mathcal{O}(1)$; we list it as a separate row (L5b) in Table~\ref{tab:taxonomy} rather than folding it into the $\mathcal{O}(1)$ suffix check.
\end{itemize}

\subsection{Cooperative Cancellation and Time-Budget Checkpointing (L6)}
\label{sec:l6-cancellation}
Execution safety and user control are provided by Layer L6 via a single atomic flag (\texttt{std::atomic<bool> stop\_flag}) and companion \texttt{solution\_found} flag. The depth-first search polls this state every $4096$ evaluated states against a configured wall-clock deadline (\texttt{time\_limit\_ms}), setting \texttt{status = UnknownTimeout} and returning gracefully if exceeded. The GUI (\texttt{dumbsspGui.cpp}) runs the solver on a dedicated background \texttt{std::thread} to preserve interface responsiveness, marshalling results back via \texttt{PostMessage(WM\_APP\_SOLVE\_FINISHED)} under a \texttt{std::mutex}. The CLI front-end (\texttt{dumbsspCli.cpp}) is fully synchronous and single-threaded, ensuring deterministic reproduction and zero thread-synchronization jitter.

\section{Bounded-Degree Zero-Sum Swap Extractor}
\label{sec:swap_manifold}

\subsection{Theoretical Foundations}
Let $S_0\subseteq\{1,\dots,n\}$ be a seed witness with $\sum_{i\in S_0}a_i=T_{\text{eff}}$, $U\subset S_0$ a removed subset, and $V\subset(\{1,\dots,n\}\setminus S_0)$ an introduced subset, giving $S' = (S_0\setminus U)\cup V$.

\begin{theorem}[Zero-Sum Delta Invariance]
\label{thm:zero_sum}
$S'$ is a valid exact solution iff $\Delta(U,V) = \sum_{v\in V}a_v - \sum_{u\in U}a_u = 0$.
\end{theorem}
\begin{proof}
$\sum_{j\in S'}a_j = \big(\sum_{j\in S_0}a_j - \sum_{u\in U}a_u\big) + \sum_{v\in V}a_v = T_{\text{eff}} + \Delta(U,V)$, which equals $T_{\text{eff}}$ iff $\Delta(U,V)=0$.
\end{proof}

\begin{theorem}[Even Hamming Distance]
\label{thm:hamming}
If $|U|=|V|=p$, the Hamming distance between the indicator vectors of $S_0$ and $S'$ is $d_H = |U|+|V| = 2p$.
\end{theorem}

\subsection{Bounded-Degree ($p \le 4$) Perturbation Extraction}
The Zero-Sum Swap Extractor (\texttt{ZeroSumSwapExtractor::extract}) generalizes local perturbation search by exhaustively exploring symmetric and asymmetric element substitutions up to degree $p \le 4$ on both the internal witness set $S_{\text{in}} = S_0$ and the external multiset $S_{\text{out}} = A \setminus S_0$. Concretely:
\begin{enumerate}
    \item Enumerate all $U\subset S_0$ with $1\le|U|\le4$: generating $\Theta(k^4)$ candidate subsets (dominated by $\binom{k}{4}$), with sums stored and sorted in $\Theta(k^4\log k)$ time.
    \item Enumerate all $V\subset S_{\text{out}}$ with $1\le|V|\le4$ ($\Theta((n-k)^4)$ subsets); for each candidate $V$, binary-search (\texttt{std::equal\_range}) the sorted $U$-list for $\sum V = \sum U$ (satisfying $\Delta(U,V)=0$, Theorem~\ref{thm:zero_sum}) in $\mathcal{O}(\log k^4)=\mathcal{O}(\log k)$ comparisons, materializing a valid witness for each match up to \texttt{max\_solutions}.
\end{enumerate}
The total extraction complexity is
\begin{equation}
    \mathcal{O}\big(k^4\log k \;+\; (n-k)^4\log k\big),
\end{equation}
which constitutes a \emph{fixed-degree polynomial} procedure in class P.

\section{Reproducible Experimental Evaluation}
\label{sec:experiments}

The experimental evaluation is based on a fully reproducible $20$-instance self-test transcript, archived in \texttt{benchmark.txt}, covering $N=5,10,15,\dots,100$ in steps of $5$, each a randomly generated flat instance with a trillion-scale target, solved in \texttt{FindOne} mode and independently re-verified by Layer L7. Every row in Table~\ref{tab:selftest} is reproduced verbatim from that transcript, including the exact CLI invocation, and can be regenerated by any reader who compiles \texttt{dumbsspCli.cpp} against \texttt{dumbsspCore.hpp}.

\subsection{Hardware Testbed \& Benchmark Environment}
All benchmark instances and timing results reported in Table~\ref{tab:selftest} were measured on a commodity low-power entry-level hardware platform with the following system configuration:
\begin{itemize}
    \item \textbf{System Model:} Acer Aspire A314-35 (System BIOS v1.09 UEFI).
    \item \textbf{Processor (CPU):} Intel(R) Celeron(R) N5100 @ 1.10\,GHz (4 physical cores, 4 threads, 4\,MB L3 cache, Jasper Lake 10\,nm, ultra-low-power 6\,W TDP).
    \item \textbf{System Memory (RAM):} 8192\,MB (8.0\,GB) DDR4 RAM (7998\,MB OS visible).
    \item \textbf{Operating System:} Microsoft Windows 11 Pro 64-bit (10.0, Build 26200).
    \item \textbf{Compilation \& Execution Model:} MinGW-w64 GCC (\texttt{-O3 -std=c++17}), executing in a strictly \textbf{single-threaded}, synchronous console mode via \texttt{dumbsspCli.exe}.
\end{itemize}
The exact problem instances, target numbers, and execution logs are archived in \texttt{benchmark.txt}.

\subsection{Reproduction Instructions}
\label{sec:repro}
We recommend reproducing all quantitative claims in this paper via the \textbf{command-line interface}, \texttt{dumbsspCli.exe} (built from \texttt{dumbsspCli.cpp} and \texttt{dumbsspCore.hpp}), rather than via the Win32 GUI (\texttt{dumbsspGui.cpp}). The CLI is fully synchronous, single-threaded, and prints the exact statistics used below. General invocation:
\begin{verbatim}
dumbsspCli.exe "a_1, a_2, ..., a_N" T [findone|findall|count|decision]
\end{verbatim}
For example, the $N=20$ row of Table~\ref{tab:selftest} is regenerated by:
\begin{verbatim}
dumbsspCli.exe "879273641953, 844911863763, 148075117490, ...,
998161222768" 3972117833710 findone
\end{verbatim}
(full element lists are given in \texttt{benchmark.txt}). Running with no arguments executes the four built-in regression checks embedded in \texttt{main()} (Test 1--4 over a fixed 32/33-element dataset) as a smoke test.

\begin{table}[h]
\centering
\caption{Reproducible Self-Test Transcript: \texttt{FindOne} on Random Trillion-Scale Instances (source: \texttt{benchmark.txt})}
\label{tab:selftest}
\renewcommand{\arraystretch}{1.1}
\begin{tabular}{cccc}
\hline
$N$ & Witness Size $k$ & Wall-Clock Time & L7 Verifier \\
\hline
5   & 2  & $0.028$\,ms  & OK \\
10  & 3  & $0.053$\,ms  & OK \\
15  & 5  & $2.210$\,ms  & OK \\
20  & 6  & $83.889$\,ms & OK \\
25  & 8  & $250.305$\,ms & OK \\
30  & 10 & $395.859$\,ms & OK \\
35  & 11 & $317.753$\,ms & OK \\
40  & 13 & $477.693$\,ms & OK \\
45  & 15 & $207.076$\,ms & OK \\
50  & 15 & $919.937$\,ms & OK \\
55  & 21 & $4.208$\,s   & OK \\
60  & 20 & $4.219$\,s   & OK \\
65  & 19 & $15.752$\,s  & OK \\
70  & 24 & $2.953$\,s   & OK \\
75  & 23 & $7.142$\,s   & OK \\
80  & 20 & $1.397$\,s   & OK \\
85  & 25 & $26.898$\,s  & OK \\
90  & 27 & $5.297$\,s   & OK \\
95  & 27 & $4.846$\,s   & OK \\
100 & 28 & $17.800$\,s  & OK \\
\hline
\end{tabular}
\end{table}

\subsection{Discussion}
All $20/20$ instances were certified satisfiable, and every returned witness passed independent L7 re-verification (unique indices, values matching the original array, exact sum equality). Two observations follow directly from the data and from the router logic in Algorithm~\ref{alg:router}:
\begin{enumerate}
    \item \textbf{All rows use L4, not L3.} Every target in the transcript is trillion-scale, far above the $1.5\times10^7$ Bitset DP threshold, so every instance is routed to the Hybrid Tail-Table engine; the transcript does not exercise L3 at all.
    \item \textbf{Runtime is not monotone in $N$.} $N=65$ ($15.752$\,s) is slower than both $N=60$ ($4.219$\,s) and $N=70$ ($2.953$\,s); $N=85$ ($26.898$\,s) is the slowest point in the whole table despite $N=90$--$100$ running faster. This is consistent with the theory of Section~\ref{sec:preliminaries}: for random flat instances, DFS cost is governed by how tightly L5a/L5b prune the $2^{\,n-m}$-node search space, which depends on the cardinality window width and local structure of each randomly drawn instance, not on $N$ alone. We therefore avoid fitting a single closed-form ``time vs.\ $N$'' curve to this data, and report the raw transcript instead.
\end{enumerate}
Peak RAM (\texttt{GetProcessMemoryInfo}) was not logged in this transcript, so we do not restate the earlier draft's per-row MB figures; the \emph{analytical} memory bound of Section~\ref{sec:hardware_engineering} (a $16.00$\,MB tail table at the default $m=20$, plus $\mathcal{O}(n)$ auxiliary suffix/path arrays) still applies and is independent of $T$ whenever the L4 path is taken.

\subsection{Qualitative Comparison Against Classical Baselines}
Table~\ref{tab:comparison} restates well-known asymptotic behavior of classical paradigms from the literature, and describes the proposed solver by its \emph{analytical} complexity (Table~\ref{tab:taxonomy}) and by the empirical range actually observed in Table~\ref{tab:selftest}, rather than by unlogged figures.

\begin{table*}[t]
\centering
\caption{Qualitative Paradigm Comparison}
\label{tab:comparison}
\renewcommand{\arraystretch}{1.3}
\footnotesize
\resizebox{\textwidth}{!}{%
\begin{tabular}{llllll}
\hline
\textbf{Paradigm} & \textbf{Time} & \textbf{Memory} & \textbf{Behavior, $N\ge50$} & \textbf{Behavior, $T>10^{12}$} & \textbf{Multi-Solution} \\
\hline
Horowitz--Sahni (1974) & $\mathcal{O}(2^{N/2})$ & $\mathcal{O}(2^{N/2})$ ($>1$\,GB) & Cache thrashing & Target-independent & $\mathcal{O}(K\cdot2^{N/2})$ \\
Pisinger \texttt{minknap} (1997) & Pseudo-P. & $\mathcal{O}(N\cdot T)$ & Fast if $T$ small & Fails/OOM if $T>10^9$ & Single witness \\
Standard DP Bitset & $\mathcal{O}(N\cdot T)$ & $\mathcal{O}(T)$ & Fast if $T$ small & Fails/OOM & Exhaustive backtrack \\
\textbf{This work (L3/L4 router)} & Adaptive (P / pseudo-P / $2^{n-m}$) & $\le16$\,MB table $+\,\mathcal{O}(n)$ (or $\mathcal{O}(T)$ if L3) & $0.028$\,ms--$26.9$\,s (Table~\ref{tab:selftest}) & Same L4 path, table size independent of $T$ & L8: $\mathcal{O}(k^4\log k{+}(n{-}k)^4\log k)$, capped at $5000$ \\
\hline
\end{tabular}%
}
\end{table*}

\section{User Interface \& Independent Verification}
\label{sec:ui_verification}

\subsection{User Interface Architecture}
The implementation provides two front-end interfaces sharing the same underlying core engine: an interactive Win32 graphical interface (\texttt{dumbsspGui.cpp}) featuring live telemetry rendering, parameter configuration, and structured-exception handling (\texttt{GlobalCrashHandler}); and a standalone console executable (\texttt{dumbsspCli.exe}). For rigorous reproducible benchmarking, headless automated testing, and timing precision, the CLI front-end is utilized to eliminate window-manager scheduling jitter and rendering latency.

\subsection{L7 Independent Black-Box Verifier}
Layer L7 (\texttt{verify\_witness\_independently}) runs after and outside the search engine. Given the original $A_{\text{raw}}$, the target $T$, and a candidate witness, it checks, using $128$-bit accumulation for the sum: (i) every index lies in $[0,N)$; (ii) no index repeats (via an \texttt{unordered\_set}); (iii) each claimed value matches $A_{\text{raw}}$ at that index; and (iv) the accumulated sum equals $T$ exactly. If Theorem~\ref{thm:complement} was applied, the witness is inverted against $A$ \emph{before} this check, so L7 always certifies against the problem as originally stated by the user.

\section{Conclusion \& Future Work}
\label{sec:conclusion}

We presented \textsc{Dumb SSP Solver}, an exact multi-paradigm solver combining memory-bounded asymmetric tail-table memoization (L4, $16.00$\,MB L3 cache-resident footprint), threshold-gated word-parallel bitset DP (L3), suffix and local-cardinality pruning (L5a/L5b), cooperative cancellation checkpointing (L6), independent black-box verification (L7), and bounded-degree ($p\le4$) zero-sum swap extraction (L8) for multi-solution harvesting. The architecture was empirically evaluated on random instances with trillion-scale targets across sizes $N=5$ to $N=100$ on commodity laptop hardware, achieving exact witness identification under bounded memory ($\le 21$\,MB).

From a practical perspective, the observed ability to evaluate subset sum instances with trillion-scale targets within seconds under a bounded memory footprint on commodity hardware suggests potential utility across several computational domains. In cryptanalysis and security, such exact search mechanisms may serve as exploratory tools for evaluating knapsack-based cryptosystems, calibrating hardness assumptions in lattice-based cryptography, or facilitating witness generation in zero-knowledge proof protocols. In systems and applied optimization, the approach potentially aids in exact bin packing for cloud resource allocation, financial transaction reconciliation, and zero-scrap cutting stock problems.

Key directions for future work include:
\begin{itemize}
    \item \textbf{Root-branch parallelization:} implementing multi-threaded work-stealing for the L4 depth-first search tree across multi-core CPU architectures.
    \item \textbf{Lattice-reduction integration:} incorporating LLL/BKZ lattice basis reduction to accelerate exact solving in low-density regimes ($d < 0.9408$).
    \item \textbf{Gray-code tail-table synthesis:} generating tail-table sums via incremental single-bit Gray-code transitions to reduce table initialization complexity to $\mathcal{O}(2^m)$.
    \item \textbf{Automated combinatorial scaling:} integrating direct $2^Z$ zero-element expansion within headless batch pipelines.
\end{itemize}

\bibliographystyle{IEEEtran}
\bibliography{references}

\end{document}
